\section{One Compilation, and an Operation That Keeps Compiling}
\label{sec:ch16-one-compilation}
\index{operation cache!what it holds|(}%

Two questions about the operation cache were left open in
chapter~\ref{ch:the-life-of-a-request}. The first was what happens when several
requests for the same uncached operation arrive together, which that chapter read
in the source and never exercised. The second nobody thought to ask, because the
word \enquote{compiled} carries an assumption with it.

\subsection*{A burst compiles once}

\index{operation cache!single-flight compilation}%
The mechanism, in \code{OperationCacheMiddleware.cs}, is a table of compilations
in progress:

\begin{minted}[fontsize=\footnotesize,breaklines]{csharp}
private readonly ConcurrentDictionary<string, Lazy<TaskCompletionSource<Operation>>> _inFlightOperations =
    new(StringComparer.Ordinal);
\end{minted}

and an election written with unusually candid comments:

\begin{minted}[fontsize=\footnotesize,breaklines]{csharp}
// No operation is cached and no compilation is in progress.
// Use a Lazy<TCS> so that under burst conditions only one TCS is materialized
// even if multiple requests race through GetOrAdd concurrently.
inFlightOperation = new Lazy<TaskCompletionSource<Operation>>(
    static () => new TaskCompletionSource<Operation>(
        TaskCreationOptions.RunContinuationsAsynchronously));
var cachedInFlightOperation = _inFlightOperations.GetOrAdd(operationId, inFlightOperation);

if (ReferenceEquals(cachedInFlightOperation, inFlightOperation))
{
    // We won the race! This request is the single-flight leader
    // responsible for compiling and signaling all followers.
    isSingleFlightLeader = true;
    context.Features.Set(inFlightOperation.Value);
}
\end{minted}

\index{operation cache!where the compile actually happens}%
Chapter~\ref{ch:the-life-of-a-request} quoted that file and left an impression I
have to correct, because it is mine. The cache middleware does not compile
anything at all. The winner stores its completion source as a request feature
and falls straight through, and the compile happens one position further down,
in \code{OperationResolverMiddleware}, which signals the waiting requests the
instant compilation returns:

\begin{minted}[fontsize=\footnotesize,breaklines]{csharp}
operation = _operationPlanner.Compile(...);

context.SetOperation(operation);
inFlightOperation?.TrySetResult(operation);
\end{minted}

The difference is not academic. What is coalesced is compilation, not the
request: the followers are released while the leader is still executing its own
query, which is what the middleware's name promises and not what I had pictured.

\index{operation cache!the leader that throws}%
The failure path is the half nobody asked about and the half I would want to
know. If the leader throws while compiling, it calls \code{TrySetException} on
the way out, so every follower faults with the leader's exception at its own
await rather than hanging or silently retrying. The in-flight entry is removed
in a \code{finally} either way, so the next request to arrive becomes a fresh
leader and tries again. There is a second guard for the case where the pipeline
somehow produces no operation and does not throw:

\begin{minted}[fontsize=\footnotesize,breaklines]{csharp}
// The pipeline completed without producing an operation and without
// throwing. Signal followers so they do not hang indefinitely.
\end{minted}

And there is no timeout in any of this. A follower waits on its own
\code{RequestAborted}, which by this point is the token
\code{TimeoutMiddleware} built from the client's abort signal and the execution
timeout, thirty seconds by default. Each follower's clock is its own and starts
when its own request entered the pipeline, so a follower can time out while the
leader is still compiling, and that affects nobody else.

\index{operation cache!measured under a burst}%
Chapter~\ref{ch:the-life-of-a-request} could not check any of this, so
\code{samples/executor-internals} does. Thirty-two concurrent requests for a
document the executor has never seen:

\begin{minted}[fontsize=\footnotesize,breaklines]{text}
== single-flight-compiles-once
   32 concurrent identical first-time requests -> 1 compilation(s)
   PASS
\end{minted}

One compilation, five runs in a row. The comments describe what the code does.

\subsection*{And then it goes on compiling}

\index{operation cache!the operation is not finished}%
Here is the assumption the word \enquote{compiled} smuggles in. I had read the
operation cache as holding finished artefacts: parse once, validate once,
compile once, and thereafter the entry is a constant that many requests read
concurrently and nobody writes.

That is not what is in there. Section~\ref{sec:ch16-what-compiling-produces}
established that \code{Compile} builds the root selection set and nothing below
it. Everything below it is built in \code{Operation.cs}, on demand, at execution
time:

\begin{minted}[fontsize=\footnotesize,breaklines]{csharp}
var key = (selection.Id, typeContext.Name);

if (!_selectionSets.TryGetValue(key, out var selectionSet))
{
    lock (_sync)
    {
        if (!_selectionSets.TryGetValue(key, out selectionSet))
        {
            selectionSet =
                _compiler.CompileSelectionSet(
                    this,
                    selection,
                    objectType,
                    _includeConditions,
                    _deferConditions,
                    ref _elementsById,
                    ref _lastId);
            _selectionSets.TryAdd(key, selectionSet);
        }
    }
}
\end{minted}

\index{selection sets!keyed by field and concrete type}%
Three things in that block are worth separating. The key is a pair, a selection
identifier and a \emph{concrete object type name}, so one field selected once in
the document compiles as many selection sets as there are concrete types that
actually come back through it. The lookup is a double-checked lock rather than a
plain read, because several in-flight requests share this object. And
\code{\_elementsById} and \code{\_lastId} go in by reference and are mutated in
place: the array of compiled elements grows and the identifier counter advances,
inside the lock, on a cached object.

So a cached operation is a partially compiled one, and what it finishes
compiling is decided by the shapes your resolvers return.
Figure~\ref{fig:ch16-lazy-compilation} is the two halves on one picture.

\begin{figure}[htbp]
  \centering
  % What Compile builds eagerly, and what execution builds later. Referenced from
% chapter 16. Bare tikzpicture; the figure environment, caption and label live
% at the call site.
%
% The point of the picture is that the dividing line runs through the middle of
% one object. The root selection set is built once at position nine; every
% selection set below an abstract field is built the first time a resolver
% actually returns that concrete type, and written into the same cached
% Operation that other requests are reading. The dashed boxes are the ones that
% do not exist yet when the operation is called compiled.
\begin{tikzpicture}[
    font=\small,
    box/.style={draw, rounded corners=2pt, minimum width=24mm,
                minimum height=7mm, align=center, font=\scriptsize},
    later/.style={draw, dashed, rounded corners=2pt, minimum width=22mm,
                  minimum height=7mm, align=center, font=\scriptsize},
    op/.style={draw, thick, rounded corners=2pt, minimum width=32mm,
               minimum height=9mm, align=center, font=\scriptsize,
               fill=black!8},
    flow/.style={-{Stealth[length=2mm]}},
    lazy/.style={-{Stealth[length=2mm]}, dashed},
    note/.style={font=\scriptsize, gray, align=center, inner sep=1pt},
    hdr/.style={font=\scriptsize\itshape, align=center}
  ]

  % -- the eager half, at position nine --------------------------------------
  \node[hdr] at (2.2,2.4) {position nine, once};
  \node[op]  (op)   at (2.2,1.5) {the compiled \code{Operation}};
  \node[box] (root) at (2.2,0.1) {root selection set};

  \draw[flow] (op) -- (root);

  % -- the lazy half, at execution time --------------------------------------
  \node[hdr] at (8.6,2.4) {later, on first sight of a type};
  \node[box] (field) at (8.6,1.5) {the \code{node} selection};

  \node[later] (t1) at (6.4,0.1)  {for \code{Product}};
  \node[later] (t2) at (10.8,0.1) {for \code{Order}};
  \node[later] (t3) at (8.6,-1.1) {for \code{Review}};

  \draw[flow] (root) -- (field);
  \draw[lazy] (field) -- (t1);
  \draw[lazy] (field) -- (t2);
  \draw[lazy] (field) -- (t3);

  \node[note, anchor=north] at (8.6,-1.5)
    {each compiled the first time a resolver returns that type,\\
     then memoised against the pair (selection id, type name)};

  % -- and what it costs the object ------------------------------------------
  \node[note, anchor=north] at (5.4,-2.6)
    {The memo table, the element array and the id counter all live on the cached
     \code{Operation},\\
     and are written under a lock while other requests read it. Both caches
     report a hit throughout.};

\end{tikzpicture}

  \caption{What \code{Compile} builds, and what the first request to see a new
  concrete type builds. The eager half runs once, at position nine, and produces
  the root selection set alone. Everything under an abstract field is compiled
  the first time a resolver actually returns that type, memoised against the
  pair of a selection identifier and a type name, and written into an object
  that other requests are reading at the same time. Both caches report a hit
  throughout, because both of them did.}
  \label{fig:ch16-lazy-compilation}
\end{figure}

\index{abstract types!and first-use compilation}%
For most schemas this is an implementation detail with no consequence you could
observe. For this book it is worth a paragraph, because abstract types cluster
exactly where a federated graph puts its seams.
Chapter~\ref{ch:hard-modeling-problems} gave Mosaic a \code{node} field that
returns an interface implemented by four types owned by four different services,
and every one of those is a concrete type that has to be seen once before its
selection set exists. An implementation that turns up rarely pays a one-off, locked
compilation on the request that first returns it, long after the operation was
compiled and cached, and every cache counter you have will report a hit while it
happens.

How far that goes is a separate question, and I have not answered it. I did not
measure the cost and I am not going to print a latency I did not take. What the source supports is the shape: the work is real, it is
deferred to first use, it is per concrete type, and it is invisible to the two
counters chapter~\ref{ch:the-life-of-a-request} taught you to read. Whether it
is ever big enough to see on a wall clock is a question for
chapter~\ref{ch:resilience-and-performance}, which owns the load testing and can
answer it properly.

\medskip
That is what the compiler makes and where it makes it. The next question is what
it decides about each field on the way, because one of those decisions has been
quietly shaping every number this book has printed.

\index{operation cache!what it holds|)}%
